% !TEX encoding = IsoLatin2  % notwendige Zeile f"ur Mac-Benutzer (muss als Kommentar stehen); Windows-Benutzer k"onnen die Zeile l"oschen.

% LaTeX-Vorlage Version 3.1,  Juli 2011
% erstellt von Dr. Andreas Drauschke (andreas.drauschke@technikum-wien.at) und Dr. Susanne Teschl (susanne.teschl@technikum-wien.at)
% geringf"ugig adaptiert von Harald Stockinger (harald.stockinger@technikum-wien.at)


\documentclass[a4paper,bibtotoc,oneside]{scrbook}
% F"ur kurze Arbeiten w"are auch die Dokumentklasse "scrartcl" ausreichend. In diesem Fall ist "section" die h"ochste Ebene ("chapter" gibt es dann nicht).
% \documentclass[a4paper,bibtotoc,oneside]{scrartcl}


%Zum Verlinken des inhaltsverzeichniss & co
\usepackage[colorlinks=false]{hyperref}
% %hyperref setup`%
\hypersetup{pdfborder=0 0 0}
%       colorlinks=false,
%       citecolor=Violet,
% %         linkcolor=Green}
\usepackage[utf8x]{inputenc}
% deutsche Anpassungen
% \usepackage[ansinew]{inputenc}
\usepackage[T1]{fontenc}
\usepackage[ngerman]{babel}
% mathematische Symbole
\usepackage{amsmath,amssymb,amsfonts,amstext}
\usepackage{xcolor}
\usepackage{calc}
\usepackage{caption}
%tabellen

% Kopfzeilen frei gestaltbar
\usepackage{fancyhdr}
\lfoot[\fancyplain{asdf}{}]{\fancyplain{}{}}
\rfoot[\fancyplain{}{}]{\fancyplain{}{}}
\cfoot[\fancyplain{}{\footnotesize\thepage}]{\fancyplain{}{\footnotesize\thepage}}
\lhead[\fancyplain{}{\footnotesize\nouppercase\leftmark}]{\fancyplain{}{}}
\chead{}
\rhead[\fancyplain{}{}]{\fancyplain{}{\footnotesize\nouppercase\sc\leftmark}}

% Farben im Dokument m"oglich
\usepackage{color}

% Schriftart Helvetica
\usepackage{helvet}
\renewcommand{\familydefault}{cmss}

% Graphiken einbinden: hier f"ur pdflatex
\usepackage[pdftex]{graphicx}
% Um pdf einzufügen
\usepackage{pdfpages}
\usepackage{array}

% H"ohe und Breite des Textk"orpers etwas gr"osser definieren
\setlength{\textheight}{225mm}
\setlength{\textwidth}{1.05\textwidth}

% weniger Warnungen wegen "uberf"ullter Boxen
\tolerance = 9999
\sloppy

% Anpassung einiger "Uberschriften
\renewcommand\figurename{Abbildung}
\renewcommand\tablename{Tabelle}

% %footers and headers
% % \usepackage{fancyhdr}
% % \pagestyle{fancy}
% \lhead{\studiumshort}
% \chead{}
% \rhead{\fachnameshort}
% \lfoot{\teilnehmeronenachname, \teilnehmertwonachname}
% % \cfoot{erstellt am: \today}
% \rfoot{\thepage}
% \renewcommand{\headrulewidth}{0.5pt}
% \renewcommand{\footrulewidth}{0.5pt}

\begin{document}

% Kopf- und Fusszeilen initiieren
\pagestyle{fancy}

% Deckblatt:
\thispagestyle{empty}
\begin{picture}(0,0)
\color{white}\sffamily
\put(-101,-412){\includegraphics[width=1.002\paperwidth]{./picture/LPS_2011.pdf}}
\put(220,-670){\includegraphics[width=0.5\textwidth]{./picture/FHTW_Logo_4c.pdf}}
\put(-30, -20){\bfseries\huge SEMINARARBEIT}
% Titel des Studienganges einf"ugen:
\put(-30,-50){\Large im Studiengang BEL3}
% Titel der Lehrveranstaltung einf"ugen:
\put(-30,-70){\Large Lehrveranstaltung Audiotechnik}
\color{black}
% Titel der Arbeit einf"ugen:
% Die Minipage wird gesetzt, damit auch mehrzeilige Titel m"oglich werden.
\put(-32,-350){
\begin{minipage}{13cm}
\bfseries\huge D-Verstärker
\end{minipage}
}
% Name der Autorin/des Autors eingeben:
% Personenkennzeichen der Autorin/des Autors eingeben:
\put(-30,-450){\large Ausgeführt von:\ Christian Schwarzgruber}
\put(+54,-470){\large \ Alexander Rössler}
\put(-30,-490){\large Matrikelnummer: 1110254027}
\put(+63,-510){\large 1110254020}
% Name der Begutachterin/des Begutachters eingeben:
\put(-30,-550){\large Begutachter: Michael Windisch}
\put(-30,-590){\large Wien, \today} % das Datum des letzten Kompilierens wird automatisch eingesetzt
\end{picture}

\newpage

\tableofcontents\thispagestyle{empty}
\newpage

\setcounter{page}{1}

% Falls die Kapitel"uberschriften zu lang f"ur die Kopfzeile oder das Inhaltsverzeichnis sind, so erzielt man
% dort Kurzformen der Kapitelbezeichnungen mittels:
% \chapter[Kurzform]{Lange "Uberschrift}
\chapter[Recherche über Klasse D-Verstärker]{Recherche über Klasse D-Verstärker}
Es gibt unzählige Möglichkeiten einen Klasse D-Verstärker zu implementieren, doch besitzen sie allesamt die selben
Grundbauelemente, wie Schmiddtrigger und Oszillator.
Bei der Recherche sind wir auf eine gute Anleitung gestoßen die auch noch dazu sehr gut erklärt ist.\cite{cae}
%=====================================================================================================================%
\section{Zusammenfassung der Recherche}
\subsection{Funktionsweise}
\subsection{Verwendungsmöglichkeiten}

%=====================================================================================================================%
\chapter{Dreicksgenerator}


\chapter{Seminarbuch}
\begin{table}[htbp]
  \centering
  \captionsetup{margin=1pt,font=small,labelfont=bf}
  \caption{Tätigkeit}
    \begin{tabular}{| c | c| }\hline
    {\bf Arbeitstag} &{\bf Tätigkeit} \\\hline
    \hline
    1. Tag   & Recherche im Internet über D-Verstärker  \\
    2. Tag   & Dimensionierung der Dreieckgenerators und Simulation \\
    3. Tag   & Durchführung Übung 3 Signalkennwerte Messung  \\
    4. Tag   & Durchführung Übung 3 Signalkennwerte Messung  \\
    4. Tag   & Durchführung Übung 3 Signalkennwerte Messung  \\
    \hline
    \end{tabular}%
  \label{tab:addlabel}%
\end{table}%
\begin{table}[htbp]
  \centering
    \captionsetup{margin=1pt,font=small,labelfont=bf}
      \caption{Arbeitsaufwand}
      \begin{tabular}{| c | c |}\hline
      {\bf Arbeitszeit(h)} &{\bf Dokumentation(h)} \\\hline
      \hline
        2   & 2 \\
        2   & 2 \\
        4   & 7 \\
      \hline
        \textbf{12}   & \textbf{14} \\
      \hline
      \end{tabular}%
    \label{tab:addlabel}%
\end{table}%





\noindent
Hier ist ein Hyperlink auf die  \href{http://www.technikum-wien.at}{Homepage} der FH Technikum Wien. Email-Adressen k"onnen so verlinkt werden: \href{mailto:homer.simpson@springfield.com}{\texttt{homer.simpson@springfield.com}}\\

\noindent
In der Bibliothek der Fachhochschule Technikum Wien gibt es verschiedene einf"uhrende B"ucher zum Thema \glqq \LaTeX \grqq, zum Beispiel \cite{kop05}, \cite{wil06} oder \cite{mgb+05d} (deutsche Version) bzw. \cite{mgb+04e} (englische Version). Empfehlenswerte Skripten f"ur \LaTeX-Einsteiger sind z.B. \cite{mj00} und \cite{mj95}. Sie sind frei im Internet verf"ugbar.


% Literaturverzeichnis
% Das Literaturverzeichnis kann auch nach einem allf"alligen Anhang positiioniert werden (siehe "`Leitfaden f"ur Bachelor- und Diplomarbeiten"', Version 2.0, Abschnitt 2.9).

% M"oglichkeit 1: Erzeugung des Literaturverzeichnisses mit BibTeX:
% Die Quellen sind in der Datei *.bib (hier Literatur.bib) einzugeben. Danach muss diese Vorlage einmal geTeXt werden,
dann BibTeX angewendet werden und
% anschliessend nochmals zweimal geTeXt werden.
% Im Text erfolgt die Zitierung mit dem Anker-Schl"usselwort, z.B. \cite{kop05}.
\bibliographystyle{IEEEtran}
\bibliography{Literatur}

% M"oglichkeit 2: Erzeugung eines Literaturverzeichnisses ohne BibTeX:
%\begin{thebibliography}{99}
%\bibitem[kop05]{kop05}
%H.~Kopka, {\em LaTeX, Band 1: Einf"uhrung}, Pearson Studium, M"unchen, 3.~Auflage, 2005.
%\bibitem[knu98]{knu98}
%F.~Mittelbach, M.~Goossens, J.~Braams, D.~Carlisle, and Ch. Rowley, {\em The LaTeX Companion},
%Addison-Wesley, 2nd edition, 2004.
%\end{thebibliography}

% Abbildungsverzeichnis
\listoffigures
\addcontentsline{toc}{chapter}{Abbildungsverzeichnis} % f"ugt den Eintrag "Abbildungsverzeichnis" im Inhaltsverzeichnis hinzu
\newpage

% Tabellenverzeichnis
\listoftables
\addcontentsline{toc}{chapter}{Tabellenverzeichnis} % f"ugt den Eintrag "Tabellenverzeichnis" im Inhaltsverzeichnis hinzu
\newpage

% Abk"urzungsverzeichnis
% Bei Verwendung der Dokumentklasse "scrartcl" ist der Befehlt \addchap{Abk"urzungsverzeichnis} durch
% \addsec{Abk"urzungsverzeichnis} zu ersetzen
\addchap{Abk"urzungsverzeichnis}
\hspace{-17mm}\begin{tabular}{>{\raggedleft}p{0.2\linewidth} p{0.75\linewidth} p{0.1\linewidth}}
www & World Wide Web \\
URL & Uniform Resource Locator
\end{tabular}

% Anh"ange
\begin{appendix}
\chapter[Erster Anhang]{Entwicklung und Aufbau eines Klasse D-Verstärkers}
% \lhead{}
% \lhead{\textcolor{blue}{\url{www.widatec.com/}}}\label{Anhang1}
% \includepdf[pages={1-24},addtotoc={{24},{section},{2},{Entwicklung und Aufbau eines Klasse
% D-Verstärkers},{D-Verstärker}}, pagecommand={\thispagestyle{fancy}},noautoscale=true,width=1.1\textwidth,offset=0cm
% 1cm]{/home/christian/FH/3.Semester/ATK/Projekt/seminararbeit1/unterlagen/CAE.pdf}
% \lhead{\studiumshort}
%

\chapter[Zweiter Anhang]{"Uberschrift des zweiten Anhangs}

Text Text Text Text Text Text Text Text Text Text Text Text Text Text Text Text Text Text Text Text Text Text Text Text ...

\end{appendix}

\end{document}
